\chapter{Introduction}
\label{ch:introduction}


\section{Motivation}

The Second Industrial Revolution established electricity as the foundation of industrial production. Present-day transmission networks, built to deliver it, power critical infrastructure globally. Energy has therefore served as a vital component for meeting the demands of continual social and technological developments. Alongside this, there has been an increased global emphasis on combating climate change \cite{Parliamentary_Office_of_Science_and_Technology2024-kk}, especially by the United Nations. It has incorporated this into its Sustainable Development Goals, namely SDG 13, which states that we must ``take urgent action to combat climate change and its impacts`` in all actions and policies \cite{unsdg13}. Thus, it has become necessary to rapidly evolve modern power systems to integrate and handle variable renewable energy (VRE) sources within the network for progress towards decarbonisation \cite{huang2023promoting}. This large-scale transition poses significant operational challenges, requiring improved reliability and dynamic flexibility. This additionally places greater pressure on grid operators who must ensure stable operation. Undoubtedly, the intermittency of renewables such as wind power and photovoltaic solar power further significantly increases the complexity of transmission networks \cite{shadi2025explainable}, with challenges related to quality, flow, stability, and balance impacting the performance of these systems \cite{sinsel2020challenges}. 

To tackle this, real-time grid topology optimisation has been noted as a cost-effective and dynamic method for handling the reconfiguration of grids to handle uncertainties and changes within the system \cite{marzari2025explaining}. Given the combinatorial nature and high-dimensional search space involved in this task, Reinforcement Learning (RL) represents a natural solution. It has been shown to be highly effective in tackling sequential decision-making problems in a plethora of domains, including but not limited to video games, robotics, and autonomous driving \cite{ghasemi2024comprehensive}. However, integration into real-world domains remains limited by unresolved challenges within the RL domain related to compliance with physical and operational constraints, managing uncertainty, and combating partial observability \cite{rashid2024power}. Power grid operations further exacerbate these problems. 

Recent research work has brought forward a reinforcement learning benchmark environment \cite{marchesini2025rl2grid} to accurately simulate realistic power grid operations, and work has been completed to successfully evaluate reinforcement learning algorithms for achieving strong survival performance. Despite certain reinforcement learning methods encompassing the ability to achieve high performance by learning an effective policy, the black-box nature of reinforcement learning policies raises significant safety concerns \cite{dulac2021challenges}. The lack of interpretability and explainability in its decision-making process hinders its seamless integration for deployment in the real world alongside expert system operators. If its decision-making process cannot be understood or justified, grid operators would be reluctant to deploy these systems. As the complexity of the model increases, the difficulty of identifying potential failures and debugging issues while simultaneously monitoring and maintaining long-horizon goals increases in turn. Thus, transparency and trustworthiness are essential requirements for such safety-critical infrastructure. 

Explainable Reinforcement Learning (XRL) \cite{milani2022survey} is a subset of explainable artificial intelligence (XAI) that aims to address this problem, ensuring explainability and interpretability, while preserving task performance. Therefore, there is a need to focus on reinforcement learning for real-time grid control, using XRL methods to ensure decisions are interpretable and actionable for grid operators. This work explores the intrinsic branch of XRL, where policies are interpretable by design, rather than being explained after training.

\section{Objectives and Contributions}

This project investigates whether interpretable reinforcement learning methods can produce controllers that are transparent and operationally reliable in power grid environments, without sacrificing survival performance. The overarching research question that motivates this project is as follows:

\begin{quote}
    Do current state-of-the-art methods for training interpretable reinforcement learning policies preserve operational performance in power grid topology environments?
\end{quote}

Our work makes the following contributions:

\begin{enumerate}
  \item \textbf{Implementation of existing interpretable policy extraction methods.} On three grids across four leaf budgets, existing approaches are studied to map the interpretable landscape for RL in power grid operations. Methods such as decision tree extraction and policy distillation methods that aim to make black-box models intrinsically interpretable are tested, and their scalability is evaluated. By implementing and evaluating select interpretable reinforcement learning methods within the RL2Grid environment, focusing primarily on decision tree acquisition, this work compares performance and interpretability against one another within topology-based grids. For this work, a reproducible evaluation protocol is developed for the scoring of interpretable grid controllers. Methods like VIPER \cite{bastani2018verifiable}, DAgger \cite{ross2011reduction}, and one-shot distillation transfer inadequately, limited in their preservation of operational performance by both the teacher oracle and the complexity of the environment.
  \item \textbf{Gated Decision Tree Policy Optimisation (GDTPO).} To address these limitations, this work proposes GDTPO, an online algorithm that adds four modifications to DTPO \cite{vos2024optimizing}. Each modification is controlled by separate flags, so disabling all four recover DTPO exactly. On the bus14 environment, median survival is raised from 29.62\% on DTPO to 51.20\% on GDTPO. GDTPO is able to recover on average 77.8\% of its PPO black-box teacher within a 16-leaf tree, indicating that GDTPO is a promising approach for training intrinsically interpretable policies. 
  \item \textbf{Identifying the current ceiling of interpretable grid control.} By analysing the scalability of these approaches, this work identifies the challenges within this domain that hinder further progress, and addresses key takeaways from the work.
\end{enumerate}

We provide a novel framework that showcases development within the domain of interpretable grid controllers, preserving human readability and survival performance in comparison with black-box models. By aiming to push the envelope within current benchmarks, we also highlight the current barriers to adoption, and the lessons learned within this domain.

\section{Outline}

Chapter~\ref{ch:background} introduces the field of reinforcement learning, delving into explainability within RL. Additionally, this literature review goes through the primary focus of power grid systems and the work that this project builds upon, including the simulated environment used for training and testing, and popular baseline methods that produce both black-box and interpretable policies. Chapter~\ref{ch:gdtpo} discusses the issues observed with an existing protocol, and presents the modified algorithm, with reasoning behind the four modifications. Chapter~\ref{ch:Methodology} showcases the evaluation methodology chosen for this project, detailing the deviation from the existing evaluation procedure, for reasons related to reproducibility and stability. Chapter~\ref{ch:Results} reports the experiments conducted on three grids across the multitude of methods, and evaluates the results based on the findings. Based on this, Chapter~\ref{ch:lessonslearned} collates all the key takeaways that were understood through the process of this research project. Chapter~\ref{ch:Conclusion} summarises the work done, defines our contribution, and details what is required for further improvements within this domain, and finally concludes this project.